\ChapterImageStar[cap:trabajos-futuros]{Trabajos Futuros}{./images/fondo.png}\label{cap:trabajos-futuros}
\mbox{}\\

El desarrollo exitoso de este proyecto de definición de aspectos de seguridad para la infraestructura de computación en la nube del Grupo 
\GRID\ ha establecido una base sólida que abre múltiples direcciones para trabajo futuro. Las oportunidades identificadas se extienden desde 
mejoras técnicas específicas hasta iniciativas de mayor alcance que pueden potenciar significativamente el impacto de la infraestructura.

\section{Definición de Aspectos de seguridad}
\noindent

\subsection{Implementación adicional de controles}
\noindent

Las tecnologías de seguridad implementadas demostraron capacidades de escalabilidad, interoperabilidad y monitoreo continuo dentro del 
entorno prototipo construido conforme a las características de la infraestructura del \GRID. La integración entre el firewall SonicWall 
NSA 6600, el IDPS Suricata y la plataforma SIEM Wazuh permitió consolidar mecanismos de protección, correlación de eventos y supervisión 
centralizada, evidenciando la viabilidad de ampliar progresivamente las capacidades de seguridad de la arquitectura propuesta.

Bajo este contexto, la arquitectura planteada habilita la futura incorporación de controles adicionales pertenecientes a estándares 
internacionales como la \ISO/IEC 27001 y la \ISO/IEC 27017, fortaleciendo aún más la postura de seguridad del entorno \GRID.

\textbf{Controles de la \ISO/IEC 27001}: La implementación de controles adicionales relacionados con gestión de incidentes, gestión de 
activos, continuidad operativa, control de accesos y gestión de proveedores permitiría incrementar las capacidades de prevención, 
detección y respuesta ante amenazas de seguridad. La integración progresiva de estos controles contribuiría al fortalecimiento de la 
gobernanza de seguridad y a una protección más robusta de los activos de información del \GRID.

\textbf{Controles de la \ISO/IEC 27017}: La incorporación de controles orientados específicamente a entornos cloud, como gestión de 
configuración, protección de datos, segregación de entornos virtualizados y supervisión de servicios cloud, permitiría mejorar los 
mecanismos de protección sobre la infraestructura y la información procesada dentro del entorno \GRID. Asimismo, estos controles 
complementarían las capacidades de monitoreo y trazabilidad ya implementadas en la arquitectura de defensa en profundidad.

\subsection{Implementación de modelo arquitectonico \ZTA}
\noindent

La implementación de un modelo arquitectónico fundamentado en los principios de Zero Trust Architecture (\ZTA) representa una oportunidad 
estratégica para fortalecer la postura de seguridad del entorno \GRID. Este enfoque se basa en el principio de ``nunca confiar, siempre 
verificar'', estableciendo que todo acceso a recursos, servicios o componentes de la infraestructura debe ser autenticado, autorizado y 
monitoreado continuamente, independientemente de la ubicación del usuario, dispositivo o segmento de red desde donde se origine la conexión.

Bajo este contexto, la arquitectura de defensa en profundidad propuesta proporciona una base adecuada para la futura adopción de principios 
\ZTA, debido a la integración de mecanismos de monitoreo continuo, correlación de eventos, segmentación de red, control de accesos y 
supervisión de actividades implementados mediante herramientas como SonicWall NSA 6600, Suricata y Wazuh. Estas capacidades permiten 
establecer procesos de validación y análisis permanente sobre el comportamiento de usuarios, dispositivos y servicios presentes dentro de 
la infraestructura.

Asimismo, la incorporación progresiva de un modelo \ZTA\ permitiría fortalecer aspectos relacionados con autenticación multifactor, 
microsegmentación de servicios, control granular de privilegios, verificación continua de identidad y supervisión contextual de accesos, 
incrementando la capacidad de respuesta ante amenazas internas y externas. De esta manera, el entorno \GRID\ podría evolucionar hacia un 
modelo de seguridad más dinámico, adaptable y alineado con las tendencias actuales de protección para infraestructuras cloud y entornos 
académicos distribuidos.

\section{Optimización y automatización de las tecnologías}
\noindent


\textbf{Optimización y automatización de la tecnología Suricata}:  
La herramienta Suricata permite incorporar mecanismos avanzados de automatización orientados a la gestión dinámica de eventos de seguridad 
y respuesta ante incidentes. Una posible estrategia consiste en integrar plataformas de automatización como \textit{n8n}, la cual puede 
recibir las alertas generadas por Suricata, procesarlas y correlacionarlas con fuentes externas de inteligencia de amenazas, tales como 
VirusTotal u otras bases de datos de reputación. A partir de ello, se podrían automatizar acciones como el bloqueo de direcciones IP 
maliciosas, generación de notificaciones al equipo de seguridad o actualización dinámica de reglas en el firewall SonicWall NSA 6600.

\textbf{Optimización y automatización de la tecnología Wazuh}:  
La plataforma SIEM Wazuh permite fortalecer los procesos de automatización y despliegue de infraestructura mediante herramientas de 
infraestructura como código, como \textit{Terraform}. Bajo este enfoque, Terraform podría utilizarse para automatizar el aprovisionamiento 
y configuración de componentes críticos de Wazuh, incluyendo servidores, indexadores, almacenamiento de registros y nodos asociados a la 
arquitectura de monitoreo, garantizando consistencia, escalabilidad y trazabilidad en futuras implementaciones del entorno \GRID.

De igual manera, Wazuh incorpora capacidades nativas de \textit{Active Response}, integración mediante APIs REST y compatibilidad con 
plataformas externas de orquestación, permitiendo automatizar tareas como generación de alertas, ejecución de respuestas activas, 
actualización de reglas y correlación avanzada de eventos. Estas capacidades pueden complementarse con herramientas como \textit{n8n}, 
posibilitando la automatización de flujos de trabajo relacionados con creación de tickets, envío de notificaciones y gestión centralizada 
de incidentes de seguridad dentro de la infraestructura.

\section{Investigación y desarrollo}
\noindent

\subsection{Estudios de rendimiento}
\noindent

\textbf{Benchmarking de herramientas de seguridad}:  
Realización de estudios comparativos orientados a evaluar el rendimiento de tecnologías SIEM, IDPS y firewalls en entornos de computación 
en la nube académicos, considerando métricas como consumo de recursos, capacidad de correlación de eventos, tiempo de respuesta y precisión 
en la detección de amenazas.

\subsection{Nuevas aplicaciones}
\noindent

\textbf{Integración de inteligencia artificial aplicada a ciberseguridad}:  
Exploración de mecanismos de análisis basados en aprendizaje automático para la identificación de comportamientos anómalos, clasificación 
automática de incidentes y generación de alertas predictivas dentro del entorno \GRID.

\textbf{Extensión hacia entornos académicos distribuidos}:  
Evaluación de la arquitectura propuesta en otros escenarios institucionales o laboratorios académicos con necesidades similares de 
monitoreo, correlación de eventos y protección de activos tecnológicos, permitiendo validar la escalabilidad y adaptabilidad del modelo de 
defensa en profundidad planteado.

\section{Transferencia de conocimiento}
\noindent

\subsection{Documentación y capacitación}
\noindent

\textbf{Recursos educativos}:  
Desarrollo de documentación técnica, guías de implementación, manuales de configuración y material didáctico orientado a facilitar la 
comprensión, despliegue y administración de la arquitectura de defensa en profundidad propuesta. Asimismo, la generación de tutoriales 
prácticos y casos de estudio permitiría fortalecer los procesos de aprendizaje relacionados con monitoreo de seguridad, correlación de 
eventos y protección de infraestructuras cloud dentro del entorno académico.

\section{Sustentabilidad y gobierno}
\noindent

\subsection{Modelo de sostenibilidad}
\noindent

\textbf{Plan de sostenibilidad financiera}: Desarrollo de un modelo que garantice el financiamiento continuo para operación, mantenimiento 
y crecimiento de la infraestructura a largo plazo.

\textbf{Políticas de uso}: Establecimiento de políticas claras que regulen el uso de recursos, prioridades de acceso y procedimientos para 
resolución de conflictos.

\subsection{Mejora continua}
\noindent

\textbf{Ciclo de evaluación}:  
Implementación de un proceso continuo de evaluación y mejora de la arquitectura de defensa en profundidad, considerando métricas de 
rendimiento, capacidad de detección, consumo de recursos, tiempos de respuesta y retroalimentación derivada de las actividades de monitoreo
y gestión de incidentes. Este proceso permitiría identificar oportunidades de optimización sobre las configuraciones, reglas y mecanismos 
de correlación implementados dentro del entorno \GRID.

\textbf{Adaptación tecnológica}:  
Establecimiento de mecanismos orientados a evaluar, integrar y validar nuevas tecnologías emergentes relacionadas con ciberseguridad, 
monitoreo de infraestructuras cloud, automatización y análisis de eventos de seguridad. La incorporación progresiva de nuevas herramientas, 
modelos de inteligencia artificial, plataformas SOAR o mecanismos de automatización permitiría fortalecer continuamente las capacidades de 
supervisión, protección y respuesta de la arquitectura propuesta.

\section{Impacto a largo plazo}
\noindent

Las líneas de trabajo futuro identificadas poseen el potencial de fortalecer progresivamente la postura de seguridad y las capacidades de 
supervisión tecnológica dentro del entorno \GRID, estableciendo una arquitectura de referencia aplicable a otros escenarios académicos e 
institucionales con necesidades similares de protección en infraestructuras cloud. La evolución gradual de mecanismos de monitoreo, 
correlación de eventos, automatización y respuesta ante incidentes permitiría consolidar un entorno más resiliente, escalable y alineado 
con estándares internacionales de seguridad de la información.

Asimismo, la integración de nuevas tecnologías orientadas a automatización, inteligencia artificial, orquestación de seguridad y modelos 
Zero Trust puede contribuir significativamente al fortalecimiento de los procesos de gobernanza, gestión de incidentes y protección de 
activos de información dentro de la Universidad del Quindío. Esto permitiría ampliar las capacidades institucionales relacionadas con 
ciberseguridad, favoreciendo el desarrollo de actividades de docencia, investigación y extensión soportadas sobre infraestructuras 
tecnológicas más seguras y supervisadas.

Finalmente, la sostenibilidad y evolución de los aspectos de seguridad propuestos dependerán de la capacidad de mantener procesos de 
mejora continua, actualización tecnológica y adaptación frente a nuevas amenazas y requerimientos institucionales. Bajo este contexto, el 
trabajo desarrollado proporciona una base técnica y metodológica que puede ser replicada y extendida en futuros proyectos orientados al 
fortalecimiento de la seguridad en entornos académicos y de computación en la nube.
